%
% connection_lease_reaping_self_reap_buggy.tex
%
% FIDELITY CONTROL for H3 (self-reap-without-renew), not the design.
% Changes only Claim: it omits the self-renewal step and the c2 /= c?
% exclusion from the sweep it performs before granting the write, so the
% caller is swept by its own sweep exactly as HubDisplay._reap_and_release
% sweeps every session -- caller included -- because HubDisplay.apply
% never calls anything that renews the caller's own lease before it does.
% See docs/connection_lease_reaping.tex, Section "H3 ---
% self-reap-without-renew", for the full trace and the goal that must
% return FOUND against this variant and NOT found against the fixed spec.
%
% Type-check:  fuzz docs/connection_lease_reaping_self_reap_buggy.tex
% Model-check:
%   probcli docs/connection_lease_reaping_self_reap_buggy.tex \
%       -model_check -nodead -goal "..." -p DEFAULT_SETSIZE 3
%

\documentclass[a4paper,10pt,fleqn]{article}
\usepackage[margin=1in]{geometry}
\usepackage{fuzz}
\usepackage[colorlinks=true,linkcolor=blue,citecolor=blue,urlcolor=blue]{hyperref}

\begin{document}

\title{lux Connection-Lease Reaping --- Fidelity Control (Self-Reap on Write)}
\author{Buggy variant: Claim sweeps its own caller, never renewing it first}
\date{September 2026}
\maketitle

% ============================================================
\section{Introduction}
% ============================================================

This is the fidelity control for H3 in
\texttt{docs/connection\_lease\_reaping.tex}, not the design. It removes
only the self-protection \texttt{Claim} otherwise gives its own caller: no
renewal before the sweep, and no exclusion of the caller from the set the
sweep releases. \texttt{HubDisplay.apply} calls
\texttt{self.\_reap\_and\_release()} -- an unconditional sweep of
\emph{every} lapsed session under \texttt{HubClientRegistry.reap} --
before it does anything else, and never calls anything that would renew
the calling connection's own lease first. A connection that only ever
contacts the Hub through writes (an initial \texttt{show}, then only
patches) never renews outside that first call; once its lease's length
has passed, its own next write sweeps and releases itself before
performing the very write that should have kept it alive. Read the main
specification first; this document states only what differs, in
Section~\ref{sec:claim}.

% ============================================================
\section{Basic Types}
% ============================================================

\begin{zed}
  [CONN, IDENT, SCENE]
\end{zed}

Unchanged from the main specification.

% ============================================================
\section{Free Types}
% ============================================================

\begin{zed}
  TSTATE ::= tup | tdown
\end{zed}

\begin{zed}
  LSTATE ::= llive | llapsed
\end{zed}

\begin{zed}
  OWNER ::= unowned | conn \ldata CONN \rdata
\end{zed}

Unchanged from the main specification.

% ============================================================
\section{State}
% ============================================================

\begin{schema}{ConnReg}
  registered : \power CONN \\
  identity : CONN \pfun IDENT \\
  transport : CONN \pfun TSTATE \\
  lease : CONN \pfun LSTATE \\
  owner : SCENE \fun OWNER
\where
  registered \subseteq \dom identity \\
  \dom identity = \dom transport \\
  \dom transport = \dom lease
\end{schema}

Unchanged from the main specification.

% ============================================================
\section{Initialization}
% ============================================================

\begin{schema}{Init}
  ConnReg'
\where
  registered' = \emptyset \\
  identity' = \emptyset \\
  transport' = \emptyset \\
  lease' = \emptyset \\
  owner' = (\lambda s : SCENE @ unowned)
\end{schema}

Unchanged from the main specification.

% ============================================================
\section{Operations}
% ============================================================

\subsection{Claim --- BUGGY: no self-renewal, no self-exclusion from the sweep}
\label{sec:claim}

\textbf{This is the one operation that differs from the design.} The
ownership-gated write itself is unchanged -- the claim guard and the
grant are exactly as in the fixed spec. What is missing is the special
treatment the fixed \texttt{Claim} gives its own caller: \texttt{lease'}
is left unchanged instead of renewing \texttt{c?} to \texttt{llive}
first, and the sweep set drops the \(c2 \neq c?\) exclusion, so it reads
\emph{every} registered connection whose lease has lapsed --
\texttt{c?} included, exactly as
\texttt{HubClientRegistry.\_sweep\_locked} sweeps every session in
\texttt{\_sessions}, caller or not, when \texttt{HubDisplay
.\_reap\_and\_release} invokes it unconditionally at the top of
\texttt{apply}.

\begin{schema}{Claim}
  \Delta ConnReg \\
  c? : CONN \\
  s? : SCENE
\where
  c? \in registered \\
  transport~c? = tup \\
  lease' = lease \\
  registered' = registered \setminus \\
  \quad~ \{ c2 : CONN | c2 \in registered \land lease~c2 = llapsed \} \\
  (owner \oplus (\{ s : SCENE | \exists c2 : CONN @ \\
    \qquad c2 \in registered \land lease~c2 = llapsed
    \land owner~s = conn(c2) \} \\
    \qquad \cross \{ unowned \}))~s? = unowned \\
  \lor \\
  (owner \oplus (\{ s : SCENE | \exists c2 : CONN @ \\
    \qquad c2 \in registered \land lease~c2 = llapsed
    \land owner~s = conn(c2) \} \\
    \qquad \cross \{ unowned \}))~s? = conn(c?) \\
  owner' = (owner \oplus (\{ s : SCENE | \exists c2 : CONN @ \\
    \qquad c2 \in registered \land lease~c2 = llapsed
    \land owner~s = conn(c2) \} \\
    \qquad \cross \{ unowned \})) \oplus \{ s? \mapsto conn(c?) \} \\
  identity' = identity \\
  transport' = transport
\end{schema}

\subsection{Connect, Renew, TransportDies, LeaseLapse, Read, Depart --- unchanged}

Identical to the main specification: only \texttt{Claim}'s
self-protection is removed.

\begin{schema}{Connect}
  \Delta ConnReg \\
  c? : CONN \\
  i? : IDENT
\where
  registered' = registered \cup \{ c? \} \\
  identity' = identity \oplus \{ c? \mapsto i? \} \\
  transport' = transport \oplus \{ c? \mapsto tup \} \\
  lease' = lease \oplus \{ c? \mapsto llive \} \\
  owner' = owner
\end{schema}

\begin{schema}{Renew}
  \Delta ConnReg \\
  c? : CONN
\where
  c? \in registered \\
  transport~c? = tup \\
  lease' = lease \oplus \{ c? \mapsto llive \} \\
  registered' = registered \\
  identity' = identity \\
  transport' = transport \\
  owner' = owner
\end{schema}

\begin{schema}{TransportDies}
  \Delta ConnReg \\
  c? : CONN
\where
  c? \in registered \\
  transport~c? = tup \\
  transport' = transport \oplus \{ c? \mapsto tdown \} \\
  registered' = registered \\
  identity' = identity \\
  lease' = lease \\
  owner' = owner
\end{schema}

\begin{schema}{LeaseLapse}
  \Delta ConnReg \\
  c? : CONN
\where
  c? \in registered \\
  lease~c? = llive \\
  lease' = lease \oplus \{ c? \mapsto llapsed \} \\
  registered' = registered \\
  identity' = identity \\
  transport' = transport \\
  owner' = owner
\end{schema}

\begin{schema}{Read}
  \Xi ConnReg
\end{schema}

\begin{schema}{Depart}
  \Delta ConnReg \\
  c? : CONN
\where
  c? \in registered \\
  registered' = registered \setminus \{ c? \} \\
  owner' = owner \oplus \\
  \quad~ (\{ s : SCENE | owner~s = conn(c?) \} \cross \{ unowned \}) \\
  identity' = identity \\
  transport' = transport \\
  lease' = lease
\end{schema}

% ============================================================
\section{Verification: the goal that must flip}
% ============================================================

Against this variant, I5/I6's negation must be \emph{found} --- the
trace in \texttt{connection\_lease\_reaping.tex}, Section H3, reaches it
in four steps and never calls \texttt{TransportDies} at all:
\texttt{Connect}(c1, i1); \texttt{Claim}(c1, s1); \texttt{LeaseLapse}(c1);
\texttt{Claim}(c1, s1) again.

\begin{verbatim}
  # must be FOUND against this variant (NOT found against the fixed spec)
  probcli docs/connection_lease_reaping_self_reap_buggy.tex \
    -model_check -nodead \
    -goal "#s.(#c.(s:SCENE & c:CONN & owner(s) = conn(c) &
                   not(c : registered) & transport(c) = tup)))" \
    -p DEFAULT_SETSIZE 3
\end{verbatim}

The same trace also satisfies the more general I2/I5 goal, since a
transport that never died is one instance of a connection that is
absent from \texttt{registered} while still owning a scene:

\begin{verbatim}
  # must be FOUND against this variant (NOT found against the fixed spec)
  probcli docs/connection_lease_reaping_self_reap_buggy.tex \
    -model_check -nodead \
    -goal "#s.(#c.(s:SCENE & c:CONN & owner(s) = conn(c) &
                   not(c : registered)))" \
    -p DEFAULT_SETSIZE 3
\end{verbatim}

I1 must stay unaffected --- this variant does not touch
\texttt{TransportDies} or \texttt{Depart}, so the chain
\texttt{TransportDies};\texttt{Depart} still completes exactly as in the
fixed spec:

\begin{verbatim}
  # must remain FOUND
  probcli docs/connection_lease_reaping_self_reap_buggy.tex \
    -model_check -nodead \
    -goal "#c.(c:CONN & c : dom(identity) & not(c : registered) &
               transport(c) = tdown)" \
    -p DEFAULT_SETSIZE 3
\end{verbatim}

\end{document}
